\begin{center}
\hypertarget{native-060-article043}{}\pdfbookmark[2]{43.}{article-043}43.
\end{center}
Ad maiorem disquisitionum praecedentium illustrationem exemplum calculi completi pro anomalia vera et radio vectore ex tempore adiicimus, ad quem finem numeros art. 38 resumemus. Statuimus itaque $e=0{,}967\,6456\,7$, $\log q=9{,}765\,6500$, $t=63{,}54400$, unde primo deducimus constantes $\log\alpha=0{,}305\,2357$, $\log\beta=8{,}221\,7364$, $\log\gamma=0{,}002\,8755$.

\SourceLeaf{061}Hinc fit $\log\alpha t=2{,}108\,3102$, cui respondet in tabula \textsc{Barker}i valor approximatus ipsius $w=99^0\,6^{\prime}$, unde derivatur $A=0{,}02292$, et ex tabula nostra $\log B=0{,}000\,0040$. Hinc argumentum correctum quo tabulam \textsc{Barker}i intrare oportet fit $=\log\dfrac{\alpha t}{B}=2{,}108\,3062$, cui respondet $w=99^0\,6^{\prime}\,13^{\prime\prime}{,}18$; dein calculus ulterior ita se habet:
\begin{center}
\begin{tabular}{@{}ll@{}}
\makebox[18em][l]{$\log\operatorname{tang}\tfrac{1}{2}w^2$\hfill .\,.\,.\,.\,.\,.\,.\,.\,.} & $0{,}138\,5934$\\[.3em]
\makebox[18em][l]{$\log\beta$\hfill .\,.\,.\,.\,.\,.\,.\,.\,.\,.\,.\,.\,.\,.} & $8{,}221\,7364$\\[.3em]
\hline
\makebox[18em][l]{$\log A$\hfill .\,.\,.\,.\,.\,.\,.\,.\,.\,.\,.\,.\,.} & $8{,}360\,3298$\\[.3em]
$A$\hfill $=$ & $0{,}022\,9260\,8$\\[.3em]
\multicolumn{2}{@{}l@{}}{hinc $\log B$ perinde ut ante;}\\[.3em]
$C$\hfill $=$ & $0{,}000\,0242$\\[.3em]
$1-\tfrac{4}{5}A+C$\hfill $=$ & $0{,}981\,6833$\\[.3em]
$1+\tfrac{1}{5}A+C$\hfill $=$ & $1{,}004\,6094$
\end{tabular}
\end{center}
\begin{center}
\begin{tabular}{@{}ll@{}}
\makebox[20em][l]{$\log\operatorname{tang}\tfrac{1}{2}w$\hfill .\,.\,.\,.\,.\,.\,.\,.\,.\,.} & $0{,}069\,2967$\\[.3em]
\makebox[20em][l]{$\log\gamma$\hfill .\,.\,.\,.\,.\,.\,.\,.\,.\,.\,.\,.\,.\,.\,.} & $0{,}002\,8755$\\[.3em]
\makebox[20em][l]{$\tfrac{1}{2}\operatorname{Comp.}\log(1-\tfrac{4}{5}A+C)$\hfill .} & $0{,}004\,0143$\\[.3em]
\hline
\makebox[20em][l]{$\log\operatorname{tang}\tfrac{1}{2}v$\hfill .\,.\,.\,.\,.\,.\,.\,.\,.\,.\,.} & $0{,}076\,1865$\\[.3em]
$\tfrac{1}{2}v$\hfill $=$ & $50^0\quad0^{\prime}\quad0^{\prime\prime}$\\[.3em]
$v$\hfill $=$ & $100\quad0\quad0$\\[.3em]
\makebox[20em][l]{$\log q$\hfill .\,.\,.\,.\,.\,.\,.\,.\,.\,.\,.\,.\,.\,.\,.} & $9{,}765\,6500$\\[.3em]
\makebox[20em][l]{$2\operatorname{Comp.}\log\cos\tfrac{1}{2}v$\hfill .\,.\,.\,.\,.\,.} & $0{,}383\,8650$\\[.3em]
\makebox[20em][l]{$\log(1-\tfrac{4}{5}A+C)$\hfill .\,.\,.\,.\,.\,.\,.} & $9{,}991\,9714$\\[.3em]
\makebox[20em][l]{$\mathrm{C.}\log(1+\tfrac{1}{5}A+C)$\hfill .\,.\,.\,.\,.} & $9{,}998\,0028$\\[.3em]
\hline
\makebox[20em][l]{$\log r$\hfill .\,.\,.\,.\,.\,.\,.\,.\,.\,.\,.\,.\,.\,.\,.} & $0{,}139\,4892$
\end{tabular}
\end{center}
\noindent Si in hoc calculo factor $B$ omnino esset neglectus, anomalia vera errorusculo $1^{\prime\prime}$ tantum (in excessu) prodiisset affecta.
